\section{O mapa do trabalho}
\label{sec:mapa}

\paragraph{O objetivo.} Calcular a probabilidade de ruptura de um talude não saturado, sob chuva ou rebaixamento, com a
resistência e a permeabilidade como campos aleatórios. Os alvos são reproduzir a Seção 6 de Vargas Ceron et al.\ (2025)
e, depois, a chuva de Rahardjo et al.\ (2007), com o modelo acoplado do NeoPZ.

\paragraph{A cadeia.} Cada elo depende do anterior e tem um documento e um código.
\begin{itemize}
  \item a física do solo não saturado;
  \item a discretização e a sua verificação;
  \item o fator de segurança;
  \item o cenário hidráulico;
  \item a incerteza espacial;
  \item a propagação por Monte Carlo;
  \item a probabilidade de falha.
\end{itemize}

\begin{figure}[H]
\centering
\begin{tikzpicture}[
  caixa/.style={draw, rounded corners=2pt, fill=#1, text width=15.2cm, inner sep=4pt, font=\small, align=left},
  caixa/.default={gray!8}, seta/.style={-Stealth, thick, gray!70!black}, node distance=2.5mm]
  \node[caixa=blue!6] (f) {\textbf{1.\ Formulação $\bm u$--$p_w$--$p_a$}: Bishop, retenção, Darcy, Mohr--Coulomb PV2,
    Euler implícito, Newton.\\
    \emph{Documentos:} \texttt{formulacao}, \path{formulacao_verick}, \path{resumo_nao_saturado}.
    \emph{Código:} \texttt{TPZMatUnsaturatedUPwPa}, \texttt{TPZUnsatAnalysis} (Seção~\ref{sec:formulacao}).};
  \node[caixa=blue!6, below=of f] (v) {\textbf{2.\ Verificação}: ponto material, água, ar, Terzaghi, Mandel,
    Liakopoulos, talude seco contra Bishop.\\
    \emph{Documentos:} revisão bibliográfica e resumo com Liakopoulos. \emph{Código:} \texttt{Unsat*},
    \texttt{LearningLiakopoulos} (Seção~\ref{sec:verificacao}).};
  \node[caixa=green!6, below=of v] (s) {\textbf{3.\ Fator de segurança}: SRM e GIM com as pressões congeladas;
    arc-length.\\
    \emph{Documento:} \path{relatorio_fs}. \emph{Código:} \texttt{TPZUnsatAnalysis::ArcLength*},
    \path{Common/SlopeSolver} (Seção~\ref{sec:fs}).};
  \node[caixa=green!6, below=of s] (c) {\textbf{4.\ Cenário hidráulico}: chuva (Rahardjo, transiente de Richards) ou
    rebaixamento (Ceron, percolação permanente).\\
    \emph{Documentos:} \path{relatorio_chuva} e o artigo da intensidade crítica. \emph{Código:}
    \path{Common/SlopeRain}, \texttt{SlopeModel}, \texttt{SlopeRandomField} (Seção~\ref{sec:chuva}).};
  \node[caixa=orange!8, below=of c] (k) {\textbf{5.\ Campos aleatórios}: Karhunen--Loève de $c'$, $\phi'$ e $k$:
    forma fechada, Galerkin e Nyström.\\
    \emph{Documentos:} o relatório de 11/11 (\texttt{relatorio-2}) e o \texttt{nymstron.nb}. \emph{Código:}
    \texttt{KLSeparavel.h}, \path{Common/SlopeFields}, \texttt{KLCilindroMika}, \texttt{VerifyRandomFields}
    (Seção~\ref{sec:campos}).};
  \node[caixa=orange!8, below=of k] (m) {\textbf{6.\ Monte Carlo}: refinamento por realização, níveis de malha,
    campanhas no fermi.\\
    \emph{Documento:} o relatório de 11/11. \emph{Código:} \texttt{SlopeRandomFieldRefina}, \texttt{SlopeMonteCarlo}
    (Seção~\ref{sec:mc}).};
  \node[caixa=red!6, below=of m] (r) {\textbf{7.\ Confiabilidade}: $P_f$ e o seu CoV; FORM e amostragem por
    importância para $P_f$ pequenos.\\
    \emph{Documento:} \texttt{confiabilidade} (pasta \path{form_MIS}). \emph{Código:} \texttt{Confiabilidade}
    (Seção~\ref{sec:confiab}).};
  \foreach \a/\b in {f/v, v/s, s/c, c/k, k/m, m/r} \draw[seta] (\a.south) -- (\b.north);
\end{tikzpicture}
\caption{A cadeia do trabalho. Azul: o modelo; verde: a estabilidade sob um cenário; laranja: a incerteza;
  vermelho: o resultado.}
\label{fig:mapa}
\end{figure}

\paragraph{Os documentos do zip e onde estão no repositório.}
\begin{center}\small
\begin{longtable}{@{}L{5.6cm} L{2.6cm} L{7.4cm}@{}}
\toprule
Arquivo do zip & Data, páginas & No repositório \\
\midrule
\endhead
\path{formulacao/formulacao.tex} & 30/10, 25 p. & \path{docs/biot_nao_saturado/formulacao} (idêntico) \\
\path{formulacao/formulacao_verick.tex} & 01/11, 31 p. & não está \\
\path{formulacao/resumo_nao_saturado.tex} & 01/11, 13 p. & não está (versão curta) \\
\path{Liakopoulos/resumo/resumo_nao_saturado.tex} & 02/11, 24 p. & \path{docs/biot_nao_saturado/resumo} (idêntico) \\
\path{revisao_bibliografica_nao_saturado_neopz.pdf} & 31/10, 20 p. & não está \\
\path{arclength/relatorio_fs.tex} & 07/11, 34 p. & \path{docs/fs_arc_length} (idêntico) \\
\path{relatorio_form/relatorio_chuva.tex} & 14/11, 12 p. & não está, nem o código dele \\
\path{formulacao_random_fields/relatorio-2.pdf} & 11/11, 36 p. & \path{SlopeRandomFieldRefina/relatorio} \\
\path{formulacao_random_fields/nymstron.nb} & 14/11 & não está (o seu caderno do Nyström) \\
\path{form_MIS/confiabilidade.*} & 14/11, 19 p. & \path{Projects/Confiabilidade}. Na sua cópia do \texttt{.wl}, o esboço
  \texttt{gTalude} da Seção 7 está descomentado: só define a função, e chamá-la exige um executável que ainda não existe \\
\path{depricated/} & 16/10 a 30/10 & a evolução da formulação (Apêndice~\ref{sec:historia}) \\
\bottomrule
\end{longtable}
\end{center}
